\section{Selfless Service to Humanity}

\begin{quotex}
There is a fundamental unity underlying the Logos, man and world. The world is in existence because the Logos has willed it so. It is His Will that holds it together. Man strives to gain union with God; and when this union is achieved, the individual will merge in the mighty Universal Will. When this is achieved, will the individual say:

``I shall do no Action and I shall not help the world'' --- the world which is, because the Will with which he has sought union has willed it to be so?

I have thus solved the question for myself and I hold that serving the world, and thus serving His Will, is the surest way of salvation; and this way can be followed by remaining in the world and not going away from it.
\flright{\textsc{Bal Gangadhar Tilak}, \emph{The Secret of the Bhagavad Gita}}
\end{quotex}


\paragraph{Bhagavad Gita}


Tilak was well respected by both Guenon and Evola, primarily on the basis of his \emph{Arctic Home in the Vedas}. He also wrote an extended commentary on the Bhagavad Gita, which he saw as a call to action, rather than to enlightenment or devotion. In his words:

\begin{quotex}
The conclusion I have come to is that the Gita advocates the performance of action in this world even after the actor has achieved the highest union with the Supreme Deity by jnana or bhakti. This action must be done to keep the world going by the right path of evolution which the Creator has destined the world to follow. In order that the action may not bind the actor, it must be done with the aim of helping His purpose, and without any attachment to the coming result.
\end{quotex}


His point is that while \emph{jnana} (knowledge) and \emph{bhakti} (love) apply to God and Man, the third term of the triad, the World, is left out, unless action in the World is also part of the teaching. So, in opposition to most commentaries on the Gita, that claim that the jnana and bhakti lead away from involvement in the world, Tilak claims the opposite, that the way of jnana and bhakti should indeed lead toward action in the world.

Of course, such action should be selfless, not in the sentimental or altruistic sense that is common today, but rather in the sense that the action follows from the Will of God, not one's personal will residing in soul life. (We have been calling this the True Will.) Such action is detached and not attached to the fruits resulting therefrom. This means, in particular, that such action is free and arises from one's very nature and does not conform to any external system or code of morality. Selfless actions, then, may sometimes appear to be cruel and immoral from the all-too-human perspective. Thus, the executioner may be providing a more selfless service to humanity than a worker in a soup kitchen, provided he acts with the correct attitude: without resentment, vengeance, or sadistic pleasure.

\paragraph{Modes of Spiritual Realization}


In the essay \emph{Modes of Spiritual Realization}, \textbf{Frithjof Schuon} takes up this theme and clarifies some points. First of all, the three modes of knowledge, love and action are not coequal, but rather have a definite relationship to each other.

\subparagraph{Knowledge}


Knowledge, or gnosis, is knowledge of principles, and hence is ``completely independent of any doctrinal formulation.'' The \emph{jnani} contemplates transcendent realities as his natural state of consciousness. He possesses the Truth in an active manner, that is, he is capable of ``expressing it spontaneously, in an original and inspired manner.''

\subparagraph{Love}


The \emph{bhakta} ``obtains everything by means of love and Divine Grace,'' so that the doctrinal concerns of the jnani are not as important. Love and beauty is his path. Again, this is not sentimentalism. A man loves his religion, his country, his comrades, his family. He absorbs his religion, which ``thinks for him'', through its symbols, scriptures, and rites, for which he shows devotion. Love and gnosis are not two separate paths, nor is love subservient to gnosis, for they arise from the Divine Intellect. Neither is discursive, but in relation to the intellect one is active and objective, the other passive and subjective.

\subparagraph{Action}


Unlike gnosis and love, action does not have its sufficient reason within itself, but depends on the other two to give it meaning.

\begin{quotex}
The path of action refers to the Divinity's aspect of Rigor, whence the connection between this path and fear; this aspect is manifested for us by the indefinity and ineluctability of cosmic vicissitudes.
\end{quotex}


Action, then, as Rigor, is an expression of the left-hand path of the Kaballah. Schuon claims that the goal of action is liberation from the cosmic vicissitudes and the goal of gnosis is liberation from Existence itself. This is where Tilak and Schuon seem to part ways, since for the former, the achievement of gnosis is expressed in action. Yet Schuon asserts that ``action must be envisaged as the accomplishment of the dharma, or duty of state, which results from the nature of the individual and it must consequently be accomplished, not only to perfection, but also without attachment to its fruits.'' Nevertheless, Schuon rejoins Tilak when he says that action is inevitable, even for the one following the path of contemplation. Therefore, there must be a path that ``is followed in the midst of worldly occupations.''

\subparagraph{Selflessness}


Schuon list three conditions of selfless action.

\begin{itemize}
\item\textit{Necessity:} The activity is a necessity not a caprice.
\item\textit{Sanctification:} The activity is offered to God out of love for God and without rebelling against destiny (\emph{amor fati}).
\item\textit{Perfection:} The activity must be perfect as to its means, its ends, and the activity itself.
\begin{itemize}
\item The activity must be proportionate to the goal to be achieved
\item The means must also be proportionate (skillful means)
\item The result must satisfy the necessity that gave rise to it
\end{itemize}
\end{itemize}


Note the difference from ``altruism'', or a false idealism, that Schuon describes as ``stupid and criminal''. Such an idealism does not realize that some things are simply unavoidable. As for Jesus' command, ``love thy neighbor as thyself'', he did not mean more than thyself, nor did he imply not to love oneself. The logic of altruism harbours an inner contradiction. If Peter gives Paul his last morsel of bread out of altruism, then Paul will return the morsel for the same reason. Hence, they are at an impasse.


\flrightit{Posted on 2011-02-01 by Cologero}

\begin{center}* * *\end{center}

\begin{footnotesize}
\begin{sffamily}

\texttt{kadambari on 2011-02-08 at 09:48 said:}

2.Note the difference from ``altruism'' or an false idealism that Schuon describes as ``stupid and criminal''

How much damage has been done in the world in the name of self-righteous altruism... \\ Yes true altruism can only take place amongst ``equals''... just as true friendship as Aristotle notes in the section on friendship in the Ethics can only take place amongst ``equals'' (intellectually and morally speaking)... . Also genuine compassion towards others is only possible by a noble character which requires correct habits and proper upbringing in childhood... 

\hfill

\texttt{Jüng Memez von Pepefeels on 2016-02-21 at 21:44 said:}

``The conclusion I have come to is that the Gita advocates the performance of action in this world even after the actor has achieved the highest union with the Supreme Deity by jnana or bhakti.''

``Even''? Post-enlightenment there is nothing else left to do!

\end{sffamily}
\end{footnotesize}

